Perform the usual ordered depth-first Euler traversal. At the entry event of a node of even depth, output the node; at the exit event of a node of odd depth, output the node. Every node has exactly one entry and one exit event, so this parity rule outputs it exactly once.

It remains to bound silent work. A silent entry event is at odd depth. The next event is either the entry of its first child, which has even depth and is output, or, after one exhaustion report when there is no child, its own odd-depth exit, which is output. A silent exit event is at even depth. If it is the root exit, the traversal terminates. Otherwise the parent has odd depth; after at most one resumed child-iterator call, the next event is either entry of an even-depth sibling, which is output, or exit of the odd-depth parent, which is output. Thus between consecutive outputs, and between the last output and termination, there are only constantly many iterator calls and polynomial-overhead stack operations. The first event is entry of the depth-zero root and is output immediately. The delay and final-termination bounds are therefore \(O(p)\), with no bound on tree depth required.